\chapter{Reflection}
\label{ch:reflection}

\section{What Worked Well}

Structuring development as a sequence of discrete, measurable optimisation
stages (baseline, DSP MAC, and INT8 packing) was the single most valuable
decision in the project. A working baseline was established and verified on
hardware before any optimisation was attempted, and each subsequent stage
changed exactly one aspect of the design: arithmetic parallelism first, then
memory traffic. This made it possible to attribute each measured improvement
to a specific architectural change, to detect regressions immediately, and to
build the evidence-based narrative used throughout Chapter~\ref{ch:results}.

The deterministic benchmark harness was equally important. Generating all
test inputs from fixed seeds at runtime, and computing golden outputs on the
same PicoRV32 host, removed the need for off-chip reference data and ensured
that the same matrices were used across all three hardware stages. Any
functional discrepancy introduced by a hardware or software change would have
been caught by the element-by-element comparison on the next rebuild.

The roofline analysis in Chapter~\ref{ch:motivation} correctly predicted before
implementation that memory traffic, not arithmetic, would be the dominant
bottleneck. The measured $1.32\times$ gain from DSP-based MAC compared to the
$2.18\times$ gain from INT8 packing confirmed this diagnosis quantitatively. A
design effort that had focused purely on widening the MAC array would have
yielded diminishing returns; addressing the memory bottleneck directly produced
the largest single performance step.

\section{What Went Wrong and What Was Learned}

The most time-consuming debugging episode involved a stale
\texttt{firmware.mif} file. The FPGA synthesis flow uses this memory
initialisation file to preload on-chip RAM with the compiled firmware image.
When the file was not regenerated after a firmware change, the board continued
to run the previous image while the expected code had in fact been rebuilt.
Diagnosing this required noticing that on-board UART output did not match the
current source, which was non-trivial when the visible output is only a few bytes. The
lesson is that FPGA flows combining independently compiled software and
hardware artefacts require explicit dependency management in the build system;
a Makefile rule that re-triggers synthesis on \texttt{firmware.mif} changes
would have prevented the issue entirely.

A second lesson came from increasing routing pressure as parallel logic was
added. Quartus reported higher congestion around the C-tile register array
and the wide MAC buses, and timing closure failed after the INT8 packing
datapath was introduced until the fitter was switched from the default
balanced optimisation mode to performance-oriented optimisation.
Parallelism in FPGA designs carries costs beyond logic area: routing,
placement, and timing closure are interdependent, and a design that fits
comfortably in ALMs can still fail to close timing if the wiring density in
a local region exceeds the device's capacity. This tension does not appear
in RTL simulation and cannot be resolved by functional correctness alone.

\section{Project Management}

The project ran across two semesters. Semester~1 covered the literature
review, the design and planning phases, RTL development and integration,
and the initial FPGA deployment. Semester~2 covered the two remaining
optimisation stages (DSP-oriented MAC and INT8 quantisation), the hardware
benchmarking campaign, and report writing. The Gantt charts in
Appendix~\ref{app:gantt} record both the original plan and the way the
work actually unfolded.

Semester~1 schedule pressure arose almost entirely from the implementation
and integration phases, which ultimately ran roughly two weeks longer than
the planned allocation, as noted in Appendix~\ref{app:gantt}. Three
factors contributed. First, RTL verification exposed a combination of
syntax errors and functional defects in the early testbenches, each of
which required targeted debugging before the unit suites would pass.
Second, the SoC was originally specified with an interrupt-driven
completion path from the accelerator to the CPU, but the unforeseen
complexity of integrating interrupt handling with the PicoRV32 exception
model led to the mechanism being descoped in favour of a simple polling
loop on the status register. Third, the plan allocated only a four-day
slack window (22--25 December) for unforeseen errors, which was consumed
almost immediately once the integration problems began to appear and
proved far too short for a project of this scale.

Semester~2 pressure originated mainly from the benchmarking phase. The
UART peripheral bundled with the PicoRV32 reference design did not
support the cycle-accurate timing readouts required by the measurement
methodology in Chapter~\ref{ch:methodology}, which blocked progress on
the benchmark harness until a replacement was identified. Adopting the
open-source \texttt{simpleuart} core resolved this and removed the
largest single source of benchmarking delay. The stale firmware image
bug discussed in the preceding section compounded the pressure by
masquerading as functional failure for several days before the root
cause was understood. Finally, after the INT8 packing datapath was
added, Quartus routing failed to close timing under the default balanced
optimisation mode; switching the fitter to performance-oriented
optimisation closed timing but added roughly twenty minutes to every
full synthesis run, which measurably lengthened the remaining debug
iterations. Preserving a fully working, deployable design at the end of
every stage, rather than accumulating partially integrated features,
was the practice that kept the overall schedule recoverable despite
these setbacks.

The overarching lesson is that hardware project timelines are dominated
by integration and debug rather than initial design. Treating build-flow
discipline (artefact management, synthesis scripts, timing closure) as
first-class engineering work from the outset, rather than as polish at
the end, would have eliminated most of the schedule pressure experienced
in the second half of the project.
